% Verbatim excerpt from the Stacks Project, chapter `proetale.tex`
% (Pro-étale Cohomology), Section 61.6 ("Identifying local rings versus
% ind-Zariski"). Retrieved from
%   https://raw.githubusercontent.com/stacks/stacks-project/master/proetale.tex
% on 2026-06-02. Tag 097D is the macro label `lemma-construct-profinite`
% (Lemma 61.6.2 in the rendered output).
%
% The surrounding tags are included verbatim so the lemma group around
% 097D can be read in context: the preceding helper 097C
% (lemma-construct), the example 09BJ (example-construct-space) that
% applies 097D, and the lemmas 097E/097F (the two w-local-morphism
% lemmas) plus proposition 097G (proposition-maps-wich-identify-local-rings)
% that consume 097D right after. The text below is copied verbatim from
% proetale.tex, lines 846–1031.

\section{Identifying local rings versus ind-Zariski}
\label{section-connected-components}

\noindent
An ind-Zariski ring map $A \to B$ identifies local rings
(Lemma \ref{lemma-ind-zariski-implies}). The converse does not hold
(Examples, Section \ref{examples-section-not-ind-etale}).
However, it turns out that there is a kind of structure theorem for
ring maps which identify local rings in terms of ind-Zariski
ring maps, see Proposition \ref{proposition-maps-wich-identify-local-rings}.

\medskip\noindent
Let $A$ be a ring. Let $X = \Spec(A)$. The space of connected
components $\pi_0(X)$ is a profinite space by
Topology, Lemma \ref{topology-lemma-spectral-pi0}
(and Algebra, Lemma \ref{algebra-lemma-spec-spectral}).

\begin{lemma}
\label{lemma-construct}
Let $A$ be a ring. Let $X = \Spec(A)$. Let $T \subset \pi_0(X)$ be a
closed subset. There exists a surjective ind-Zariski ring map $A \to B$
such that $\Spec(B) \to \Spec(A)$ induces a homeomorphism of $\Spec(B)$
with the inverse image of $T$ in $X$.
\end{lemma}

\begin{proof}
Let $Z \subset X$ be the inverse image of $T$. Then $Z$ is the intersection
$Z = \bigcap Z_\alpha$ of the open and closed subsets of $X$ containing $Z$,
see Topology, Lemma \ref{topology-lemma-closed-union-connected-components}.
For each $\alpha$ we have $Z_\alpha = \Spec(A_\alpha)$ where
$A \to A_\alpha$ is a local isomorphism (a localization at an idempotent).
Setting $B = \colim A_\alpha$ proves the lemma.
\end{proof}

\begin{lemma}
\label{lemma-construct-profinite}
Let $A$ be a ring and let $X = \Spec(A)$. Let $T$ be a profinite space and
let $T \to \pi_0(X)$ be a continuous map. There exists an
ind-Zariski ring map $A \to B$ such that with $Y = \Spec(B)$ the diagram
$$
\xymatrix{
Y \ar[r] \ar[d] & \pi_0(Y) \ar[d] \\
X \ar[r] & \pi_0(X)
}
$$
is cartesian in the category of topological spaces and such that
$\pi_0(Y) = T$ as spaces over $\pi_0(X)$.
\end{lemma}

\begin{proof}
Namely, write $T = \lim T_i$ as the limit of an inverse system finite
discrete spaces over a directed set (see
Topology, Lemma \ref{topology-lemma-profinite}). For each $i$ let
$Z_i = \Im(T \to \pi_0(X) \times T_i)$. This is a closed subset.
Observe that $X \times T_i$ is the spectrum of $A_i = \prod_{t \in T_i} A$
and that $A \to A_i$ is a local isomorphism. By Lemma \ref{lemma-construct}
we see that $Z_i \subset \pi_0(X \times T_i) = \pi_0(X) \times T_i$
corresponds to a surjection $A_i \to B_i$ which is ind-Zariski
such that $\Spec(B_i) = X \times_{\pi_0(X)} Z_i$ as subsets of
$X \times T_i$. The transition maps $T_i \to T_{i'}$ induce maps
$Z_i \to Z_{i'}$ and $X \times_{\pi_0(X)} Z_i \to X \times_{\pi_0(X)} Z_{i'}$.
Hence ring maps $B_{i'} \to B_i$
(Lemmas \ref{lemma-local-isomorphism-fully-faithful} and
\ref{lemma-ind-zariski-implies}).
Set $B = \colim B_i$. Because $T = \lim Z_i$ we have
$X \times_{\pi_0(X)} T = \lim  X \times_{\pi_0(X)} Z_i$
and hence $Y = \Spec(B) = \lim \Spec(B_i)$
fits into the cartesian diagram
$$
\xymatrix{
Y \ar[r] \ar[d] & T \ar[d] \\
X \ar[r] & \pi_0(X)
}
$$
of topological spaces. By Lemma \ref{lemma-silly}
we conclude that $T = \pi_0(Y)$.
\end{proof}

\begin{example}
\label{example-construct-space}
Let $k$ be a field. Let $T$ be a profinite topological space.
There exists an ind-Zariski ring map $k \to A$ such that
$\Spec(A)$ is homeomorphic to $T$. Namely, just apply
Lemma \ref{lemma-construct-profinite} to $T \to \pi_0(\Spec(k)) = \{*\}$.
In fact, in this case we have
$$
A = \colim \text{Map}(T_i, k)
$$
whenever we write $T = \lim T_i$ as a filtered limit with each $T_i$ finite.
\end{example}

\begin{lemma}
\label{lemma-w-local-morphism-equal-points-stalks-is-iso}
Let $A \to B$ be ring map such that
\begin{enumerate}
\item $A \to B$ identifies local rings,
\item the topological spaces $\Spec(B)$, $\Spec(A)$ are w-local,
\item $\Spec(B) \to \Spec(A)$ is w-local, and
\item $\pi_0(\Spec(B)) \to \pi_0(\Spec(A))$ is bijective.
\end{enumerate}
Then $A \to B$ is an isomorphism
\end{lemma}

\begin{proof}
Let $X_0 \subset X = \Spec(A)$ and $Y_0 \subset Y = \Spec(B)$ be the
sets of closed points. By assumption $Y_0$ maps into $X_0$ and
the induced map $Y_0 \to X_0$ is a bijection.
As a space $\Spec(A)$ is the disjoint union of the spectra
of the local rings of $A$ at closed points.
Similarly for $B$. Hence $X \to Y$ is a bijection.
Since $A \to B$ is flat we have going down
(Algebra, Lemma \ref{algebra-lemma-flat-going-down}).
Thus Algebra, Lemma \ref{algebra-lemma-unique-prime-over-localize-below}
shows for any prime $\mathfrak q \subset B$ lying over
$\mathfrak p \subset A$ we have $B_\mathfrak q = B_\mathfrak p$.
Since $B_\mathfrak q = A_\mathfrak p$ by assumption, we
see that $A_\mathfrak p = B_\mathfrak p$ for all primes $\mathfrak p$
of $A$. Thus $A = B$ by
Algebra, Lemma \ref{algebra-lemma-characterize-zero-local}.
\end{proof}

\begin{lemma}
\label{lemma-w-local-morphism-equal-stalks-is-ind-zariski}
Let $A \to B$ be ring map such that
\begin{enumerate}
\item $A \to B$ identifies local rings,
\item the topological spaces $\Spec(B)$, $\Spec(A)$ are w-local, and
\item $\Spec(B) \to \Spec(A)$ is w-local.
\end{enumerate}
Then $A \to B$ is ind-Zariski.
\end{lemma}

\begin{proof}
Set $X = \Spec(A)$ and $Y = \Spec(B)$. Let $X_0 \subset X$ and
$Y_0 \subset Y$ be the set of closed points. Let $A \to A'$ be the ind-Zariski
morphism of affine schemes such that with $X' = \Spec(A')$ the diagram
$$
\xymatrix{
X' \ar[r] \ar[d] & \pi_0(X') \ar[d] \\
X \ar[r] & \pi_0(X)
}
$$
is cartesian in the category of topological spaces and such that
$\pi_0(X') = \pi_0(Y)$ as spaces over $\pi_0(X)$, see
Lemma \ref{lemma-construct-profinite}. By
Lemma \ref{lemma-silly} we see that $X'$ is w-local and
the set of closed points $X'_0 \subset X'$ is the inverse image of $X_0$.

\medskip\noindent
We obtain a continuous map $Y \to X'$ of underlying topological spaces
over $X$ identifying $\pi_0(Y)$ with $\pi_0(X')$. By
Lemma \ref{lemma-local-isomorphism-fully-faithful}
(and Lemma \ref{lemma-ind-zariski-implies})
this corresponds to a morphism of affine schemes $Y \to X'$
over $X$. Since $Y \to X$ maps $Y_0$ into $X_0$ we see that
$Y \to X'$ maps $Y_0$ into $X'_0$, i.e., $Y \to X'$ is w-local.
By Lemma \ref{lemma-w-local-morphism-equal-points-stalks-is-iso}
we see that $Y \cong X'$ and we win.
\end{proof}

\noindent
The following proposition is a warm up for the type of result
we will prove later.

\begin{proposition}
\label{proposition-maps-wich-identify-local-rings}
Let $A \to B$ be a ring map which identifies local rings.
Then there exists a faithfully flat, ind-Zariski ring map
$B \to B'$ such that $A \to B'$ is ind-Zariski.
\end{proposition}

\begin{proof}
Let $A \to A_w$, resp. $B \to B_w$ be the faithfully flat, ind-Zariski ring
map constructed in Lemma \ref{lemma-make-w-local} for $A$, resp.\ $B$.
Since $\Spec(B_w)$ is w-local, there exists a unique factorization
$A \to A_w \to B_w$ such that $\Spec(B_w) \to \Spec(A_w)$ is w-local
by Lemma \ref{lemma-universal}. Note that $A_w \to B_w$ identifies
local rings, see
Lemma \ref{lemma-local-isomorphism-permanence}.
By Lemma \ref{lemma-w-local-morphism-equal-stalks-is-ind-zariski}
this means $A_w \to B_w$ is ind-Zariski. Since $B \to B_w$ is
faithfully flat, ind-Zariski (Lemma \ref{lemma-make-w-local})
and the composition $A \to B \to B_w$ is ind-Zariski
(Lemma \ref{lemma-composition-ind-zariski})
the proposition is proved.
\end{proof}
